\documentclass[fontsize=15pt]{jlreq}

%プリアンブル
\usepackage{amsmath}
\usepackage{mathtools}
\pagestyle{empty}
\usepackage[top=15mm, bottom=20mm, left=20mm, right=20mm]{geometry}

\newcommand{\m}[1]{\raisebox{0.1em}{\textcircled{\scriptsize #1}}}



%本文
\begin{document}
\begin{center}
{\LARGE {振り返り}} \\
-物理\m{6}- %単元ごとに変えるか消す
\end{center}
\vspace{1em}
\begin{flushright}
名前：\underline{\hspace{5cm}}
\end{flushright}

% 問題部分
\begin{enumerate}
    \item 仕事の単位をかけ．\\
    (\hspace{3cm})
    \vspace{2em} 

    \item 次の仕事の値を求めよ(単位も書いて)．\\
    \begin{itemize}
        \item[(1)] 質量3[kg]の物体を2[m]だけ手前に引っ張った．この時，手が物体にした仕事．\\
        (\hspace{4cm})
        \item[(2)] 2[kg]の机の上に300[g]のカバンが置いてある．この状態のまま手で机を持って50[cm]だけ上に持ち上げた．このとき，机がカバンにした仕事．\\
        (\hspace{4cm})
    \end{itemize}
    \vspace{2em}

    \item 仕事を$W$[J]，物体に加えた力を$F$[N]，移動距離を$x$[m]として，この3つの関係式をかけ．\\
    (\hspace{15cm})
    \vspace{2em}

    \item 位置エネルギーを$U$[J]，物体の重量を$M$[N]，机からの高さを$x$[m]としたとき，位置エネルギーの式を表わせ．ただし，机の位置を重力による位置エネルギーの基準面とする．\\
    (\hspace{15cm})
\end{enumerate}
エクストラ：問題3と問題4のエネルギーの値が等しいとき，物体に加えた力$F$[N]を$U,M,x$の内，必要な文字を使って表せ．\\
(\hspace{15cm})

\end{document}